\documentclass[11pt,a4paper]{article}

\usepackage[utf8]{inputenc}
\usepackage[T1]{fontenc}
\usepackage{amsmath,amssymb,amsthm,mathrsfs}
\usepackage{geometry}
\geometry{margin=1in}
\usepackage{hyperref}
\hypersetup{colorlinks=true,linkcolor=blue,citecolor=blue}
\usepackage{enumitem}
\usepackage{booktabs}
\usepackage{array}
\usepackage{tabularx}
\usepackage{mdframed}
\usepackage{xcolor}
\usepackage{tocloft}
\setcounter{tocdepth}{3}
\setlength{\parindent}{0pt}
\setlength{\parskip}{0.5em}

% Theorem environments
\newtheorem{theorem}{Theorem}[section]
\newtheorem{proposition}[theorem]{Proposition}
\newtheorem{lemma}[theorem]{Lemma}
\newtheorem{corollary}[theorem]{Corollary}
\theoremstyle{definition}
\newtheorem{definition}[theorem]{Definition}
\newtheorem{remark}[theorem]{Remark}

% Boxes
\newmdenv[linecolor=black,backgroundcolor=gray!8,linewidth=0.6pt,
  innertopmargin=6pt,innerbottommargin=6pt,
  innerleftmargin=8pt,innerrightmargin=8pt]{lockedbox}

\newmdenv[linecolor=blue!60,backgroundcolor=blue!4,linewidth=0.6pt,
  innertopmargin=6pt,innerbottommargin=6pt,
  innerleftmargin=8pt,innerrightmargin=8pt]{derivbox}

\newmdenv[linecolor=red!60,backgroundcolor=red!4,linewidth=0.6pt,
  innertopmargin=6pt,innerbottommargin=6pt,
  innerleftmargin=8pt,innerrightmargin=8pt]{resultbox}

% Macros
\newcommand{\R}{\mathbb{R}}
\newcommand{\E}{\mathbb{E}}
\newcommand{\T}{^{\mathsf{T}}}
\newcommand{\norm}[1]{\left\|#1\right\|}
\newcommand{\inner}[2]{\langle #1,\,#2\rangle}
\newcommand{\Fbase}{F_{\mathrm{base}}}
\newcommand{\gX}{g_X}
\newcommand{\half}{\tfrac{1}{2}}
\newcommand{\diag}{\mathrm{diag}}
\newcommand{\tr}{\mathrm{tr}}
\newcommand{\rank}{\mathrm{rank}}
\newcommand{\mfls}{\mathrm{MFLS}}
\newcommand{\bsdt}{\mathrm{BSDT}}

\title{\textbf{Complete Derivations}\\[0.4em]
\large Every Object and Variable in the Canonical--BSDT Framework\\[0.2em]
\normalsize From First Principles to Final Formula}
\author{}
\date{}

\begin{document}
\maketitle
\tableofcontents
\newpage

%=============================================================
\section{Primitives: State, Features, Metric}
%=============================================================

\subsection{State matrix $X$}

\begin{definition}[State]
The state is a matrix
\[
X \in \R^{N \times d}
\]
where $N$ is the number of agents and $d$ is the feature dimension per agent.
Row $i$ is the state of agent $i$: $x_i \in \R^d$.
In vectorised form: $\mathrm{vec}(X) \in \R^{Nd}$.
\end{definition}

\textbf{Dimensions.}
\begin{center}
\begin{tabular}{lll}
\toprule
Symbol & Meaning & Type \\
\midrule
$N$ & number of agents & positive integer \\
$d$ & features per agent & positive integer \\
$n = Nd$ & total state dimension & positive integer \\
$X$ & full state & $\R^{N\times d}$ \\
$x_i$ & state of agent $i$ & $\R^d$ \\
\bottomrule
\end{tabular}
\end{center}

\subsection{Feature map $S$}

\begin{definition}[Feature map]
\[
S : \R^{N\times d} \to \R^k, \qquad X \mapsto S(X)
\]
$S(X)$ encodes the geometric content relevant to the control objective.
$k \geq 1$ is fixed per domain.
The Jacobian of $S$ at $X$ is
\[
J(X) = \frac{\partial S}{\partial X} \in \R^{k \times n}, \quad n = Nd.
\]
\end{definition}

\textbf{Jacobian in index notation.}
Let $S = (S^1,\dots,S^k)\T$ and $X = (X^1,\dots,X^n)\T$ (vectorised). Then:
\[
J_{ab}(X) = \frac{\partial S^a}{\partial X^b}, \quad a \in \{1,\dots,k\},\; b \in \{1,\dots,n\}.
\]

\textbf{Chain rule.} For any trajectory $X(t)$:
\[
\dot S = J(X)\dot X \in \R^k.
\]

\subsection{Metric $G$}

\begin{definition}[Feature-space metric]
\[
G \in \R^{k\times k}, \quad G = G\T \succ 0.
\]
Eigenvalue bounds: $0 < \mu_G \leq \lambda(G) \leq M_G$.
\end{definition}

\textbf{Derivation of eigenvalue bounds.}
Since $G \succ 0$, all eigenvalues are positive. Order them
$0 < \lambda_1 \leq \cdots \leq \lambda_k$.
Define:
\[
\mu_G := \lambda_1 = \lambda_{\min}(G), \qquad M_G := \lambda_k = \lambda_{\max}(G).
\]
For any $v \in \R^k$:
\[
\mu_G \norm{v}^2 \leq v\T G v \leq M_G \norm{v}^2.
\]
These bounds are tight: achieved at the eigenvectors of $G$.

\subsection{Derivation of $\mu_G$ and $M_G$ for BSDT}

Under the BSDT instantiation $G = \Sigma_0^{-1}$:
\[
\mu_G = \lambda_{\min}(\Sigma_0^{-1}) = \frac{1}{\lambda_{\max}(\Sigma_0)},
\qquad
M_G = \lambda_{\max}(\Sigma_0^{-1}) = \frac{1}{\lambda_{\min}(\Sigma_0)}.
\]

Proof: if $\Sigma_0 v = \lambda v$ then $\Sigma_0^{-1} v = \lambda^{-1} v$, so eigenvalues invert. $\square$

Conditioning:
\[
\kappa(G) = \frac{M_G}{\mu_G} = \frac{\lambda_{\max}(\Sigma_0)}{\lambda_{\min}(\Sigma_0)} = \kappa(\Sigma_0).
\]

%=============================================================
\section{Energy $E$}
%=============================================================

\subsection{Definition}

\begin{lockedbox}
\[
E(X) = S(X)\T G\, S(X) \in \R_{\geq 0}
\]
\end{lockedbox}

\subsection{Properties}

\begin{proposition}[Non-negativity and zero set]
$E(X) \geq 0$ with equality if and only if $G^{1/2}S(X) = 0$,
i.e.\ $S(X) = 0$ (since $G \succ 0$).
\end{proposition}

\begin{proof}
$E = S\T G S = \norm{G^{1/2}S}^2 \geq 0$.
Equality holds iff $G^{1/2}S = 0$ iff $S = 0$ since $G^{1/2}$ is invertible.
\end{proof}

\subsection{Time derivative of $E$}

\begin{proposition}[Energy velocity]
Along any trajectory $X(t)$:
\[
\dot E(t) = \langle \gX(X(t)),\, \dot X(t)\rangle
\]
where $\gX = \nabla_X E$ is derived in Section~\ref{sec:gX}.
\end{proposition}

\begin{proof}
By the chain rule:
\[
\dot E = \frac{d}{dt}[S\T G S] = 2S\T G \dot S = 2S\T G J\dot X = (2J\T G S)\T \dot X = \gX\T \dot X.
\]
\end{proof}

\subsection{BSDT instantiation of $E$}

With $S = \tilde X \Sigma_0^{-1/2}$ (vectorised), $G = I$:
\[
E_\bsdt = \norm{\tilde X \Sigma_0^{-1/2}}_F^2 = \tr(\tilde X \Sigma_0^{-1} \tilde X\T)
\]
where $\tilde X = X - \mathbf{1}\mu_0\T$.

With $S = B(X) = (\delta_C,\delta_G,\delta_A,\delta_T)\T$, $G = A$:
\[
E_\bsdt = B(X)\T A\, B(X) = E_\mathrm{BS}(S).
\]

%=============================================================
\section{State-Space Gradient $g_X$}
\label{sec:gX}
%=============================================================

\subsection{Definition}

\begin{lockedbox}
\[
\gX(X) = \nabla_X E(X) = 2J(X)\T G\, S(X) \in \R^n
\]
\end{lockedbox}

\subsection{Derivation from first principles}

\begin{proof}
Write $E = \sum_{a,b} G_{ab} S^a S^b$.
Differentiate with respect to $X^j$:
\[
\frac{\partial E}{\partial X^j}
= \sum_{a,b} G_{ab}\left(\frac{\partial S^a}{\partial X^j} S^b + S^a \frac{\partial S^b}{\partial X^j}\right)
= 2\sum_{a,b} G_{ab} J_{aj} S^b
= 2[J\T G S]_j.
\]
Hence $\nabla_X E = 2J\T G S$.
\end{proof}

\subsection{Matrix form}

In matrix notation with $X$ as an $N\times d$ matrix (not vectorised),
the gradient is an $N\times d$ matrix:
\[
\nabla_X E = 2J\T G S \in \R^{N\times d}
\]
where $J \in \R^{k\times Nd}$ acts on $\mathrm{vec}(\dot X)$ and the result is reshaped.

\subsection{Norm of $g_X$}

\begin{proposition}[Gradient norm]
\[
\norm{\gX}^2 = 4 S\T G J J\T G S = 4 S\T G\, \mathcal{G}\, G S
\]
where $\mathcal{G} = JJ\T \in \R^{k\times k}$ is the Gram matrix of $J$.
\end{proposition}

\begin{proof}
$\norm{\gX}^2 = (2J\T GS)\T(2J\T GS) = 4 S\T G J J\T G S$.
\end{proof}

\subsection{Coercivity bounds on $\norm{g_X}$}

\begin{proposition}[Two-sided gradient coercivity]
Under (A1)--(A2) with $\sigma_{\min}(J) \geq \sigma > 0$ and
$\norm{J}_\mathrm{op} \leq J_{\max}$:
\[
C_1 \sqrt{E} \leq \norm{\gX} \leq C_2\sqrt{E}
\]
where
\[
C_1 = \frac{2\sigma\mu_G}{\sqrt{M_G}}, \qquad C_2 = \frac{2J_{\max} M_G}{\sqrt{\mu_G}}.
\]
\end{proposition}

\begin{proof}
\textbf{Lower bound.}
\[
\norm{\gX}^2 = 4\norm{J\T GS}^2
\geq 4\sigma_{\min}(J)^2 \norm{GS}^2
\geq 4\sigma^2 \mu_G^2 \norm{S}^2.
\]
Since $E = S\T G S \leq M_G\norm{S}^2$, we have $\norm{S}^2 \geq E/M_G$.
Therefore $\norm{\gX}^2 \geq 4\sigma^2\mu_G^2 E/M_G = C_1^2 E$.

\textbf{Upper bound.}
\[
\norm{\gX} = 2\norm{J\T GS} \leq 2\norm{J}_\mathrm{op}\norm{G}_\mathrm{op}\norm{S}
\leq 2J_{\max} M_G \norm{S}.
\]
Since $E = S\T GS \geq \mu_G\norm{S}^2$, we have $\norm{S} \leq \sqrt{E/\mu_G}$.
Therefore $\norm{\gX} \leq 2J_{\max}M_G\sqrt{E/\mu_G} = C_2\sqrt{E}$. $\square$
\end{proof}

\subsection{BSDT pullback gradient $\tilde G_t$}

In the BSDT system, the gradient decomposes per channel:
\[
\gX = 2\sum_{k \in \{C,G,A,T\}} [AS]_k \frac{\partial \delta_k}{\partial X} =: \tilde G_t.
\]
This is the pullback gradient of Section~XVI.3 of the anatomy document.
The canonical formula $\gX = 2J\T GS$ recovers this exactly under the identification
$G = A$, $S = B(X)$, $J = \partial B/\partial X$.

%=============================================================
\section{MFLS: Mahalanobis Field Line Score}
%=============================================================

\subsection{BSDT definition}

\[
\mfls_\bsdt(X) = \norm{\nabla_X E_\bsdt}_F = 2\norm{\tilde X \Sigma_0^{-1}}_F
\]

\subsection{Canonical derivation}

\begin{lockedbox}
\[
\mfls(X) = \norm{\gX} = 2\norm{J\T GS}
\]
\end{lockedbox}

\begin{proof}
By definition $\mfls = \norm{\nabla_X E}_F$.
The canonical gradient is $\nabla_X E = \gX = 2J\T GS$.
Therefore $\mfls = \norm{2J\T GS} = \norm{\gX}$.
\end{proof}

\subsection{Expanded forms}

Three equivalent canonical expressions:

\begin{center}
\begin{tabular}{lll}
\toprule
Form & Expression & When to use \\
\midrule
Gradient norm & $\norm{\gX}$ & Free — already computed in canonical loop \\
Explicit & $2\norm{J\T GS}$ & When gradient vector is needed \\
Gram form & $2\sqrt{S\T G\mathcal{G}GS}$ & When Gram structure is analysed \\
\bottomrule
\end{tabular}
\end{center}

\subsection{MFLS squared}

\[
\mfls^2 = \norm{\gX}^2 = 4S\T G\mathcal{G}GS = 4S\T G J J\T G S
\]

\subsection{MFLS as Fisher information}

Under $G = I_S$ (Fisher metric):
\[
\mfls^2 = 4\tr(\Sigma_0^{-1}\tilde X\T \tilde X \Sigma_0^{-1}) = 4\tr(I_\mathrm{Fisher})
\]
MFLS squared is four times the trace of the Fisher information matrix at the current state.
High MFLS means the system is in a region of high log-likelihood curvature.

\subsection{MFLS as leading indicator (proved)}

From the energy derivative with $\Fbase = 0$:
\[
\dot E = -(1-\gamma)\norm{\gX}^2 = -(1-\gamma)\cdot\mfls^2.
\]
MFLS squared is the instantaneous energy descent rate.
A spike in MFLS precedes a drop in $E$. This is a theorem, not an empirical observation.

\subsection{Verification against BSDT}

Under $S = B(X)$, $G = \Sigma_0^{-1}$, $J = \partial B/\partial X$ with
$J_C = 2\tilde X\Sigma_0^{-1}$ (the camouflage Jacobian):
\[
\gX = 2J_C\T \Sigma_0^{-1} S = 2\tilde X\Sigma_0^{-1},
\qquad
\norm{\gX}_F = 2\norm{\tilde X\Sigma_0^{-1}}_F = \mfls_\bsdt. \;\checkmark
\]

%=============================================================
\section{Adaptive Gain $\gamma$}
%=============================================================

\subsection{Definition}

\begin{lockedbox}
\[
\gamma(X) = \frac{E(X)}{E(X) + \theta}, \qquad \theta > 0
\]
\end{lockedbox}

\subsection{Derivation from the saddle-point game}

Consider the two-player zero-sum game at fixed $X$:
\begin{itemize}
\item Controller chooses $\gamma \in [0,1)$
\item Adversary chooses $u \in \R^n$ with $\norm{u} \leq M_\adv$
\item Payoff to adversary: $J(\gamma,u) = \dot E(X;\gamma,u)$
\end{itemize}

The payoff under the canonical ODE is:
\[
J(\gamma,u) = (1-\gamma)\left[\inner{\gX}{\Fbase + u} - \norm{\gX}^2\right].
\]

The adversary maximises: by Cauchy--Schwarz, $u^* = M_\adv \gX/\norm{\gX}$.

The controller minimises:
\[
\inf_{\gamma \in [0,1)} (1-\gamma) B, \quad B = \inner{\gX}{\Fbase} + M_\adv\norm{\gX} - \norm{\gX}^2.
\]
When $B > 0$: infimum is $0$ as $\gamma \to 1^-$; boundary solution.
When $B \leq 0$: minimum at $\gamma = 0$.

The Lagrangian formulation (Theorem O.1 of v4):
\[
\mathcal{L}(u,\lambda) = \half\norm{u - \Fbase}^2 + \lambda\inner{u}{\gX} - \frac{\theta}{2}\frac{\norm{\gX}^2}{E}\lambda^2
\]
has unique saddle point with $\lambda^* = \frac{\inner{\Fbase,\gX}E}{(E+\theta)\norm{\gX}^2}$,
giving $\gamma = E/(E+\theta)$.

\subsection{Properties}

\begin{proposition}[Gain properties]
\begin{enumerate}[label=(\roman*)]
\item $\gamma \in [0,1)$ for all $X$
\item $\gamma \to 0$ as $E \to 0$: no damping near target
\item $\gamma \to 1$ as $E \to \infty$: full projection at high energy
\item $\gamma = 1/2$ when $E = \theta$: maximum-entropy operating point
\item $\gamma' = \theta/(E+\theta)^2 > 0$: strictly increasing in $E$
\end{enumerate}
\end{proposition}

\begin{proof}
(i) $E \geq 0$, $\theta > 0$, so $0 \leq E/(E+\theta) < 1$.
(ii) $\lim_{E\to0} E/(E+\theta) = 0$.
(iii) $\lim_{E\to\infty} E/(E+\theta) = 1$.
(iv) $E/(E+\theta) = 1/2 \iff E = \theta$.
(v) $d\gamma/dE = \theta/(E+\theta)^2 > 0$.
\end{proof}

\subsection{Fermi-Dirac structure}

Writing $\gamma = 1/(1 + \theta/E)$: this is the Fermi-Dirac occupation function
with inverse temperature $\beta = 1/\theta$ and energy $E$.
The thermodynamic dictionary (v4 §N):
\begin{center}
\begin{tabular}{ll}
\toprule
Canonical object & Thermodynamic analogue \\
\midrule
$\theta$ & Temperature \\
$\gamma$ & Occupation number \\
$1-\gamma$ & Carnot efficiency \\
$E = \theta$ & Maximum-entropy operating point \\
$\theta \to 0$ & Zero temperature (full intervention) \\
$\theta \to \infty$ & Infinite temperature (no intervention) \\
\bottomrule
\end{tabular}
\end{center}

\subsection{Formula for $\theta$}

\begin{lockedbox}
\[
\theta = \mathrm{median}_{t \in \text{calm}}\{E(t)\}
= \mathrm{median}_{t \in \text{calm}}\left\{\norm{\tilde X_t \Sigma_0^{-1/2}}_F^2\right\}
\]
\end{lockedbox}

\textbf{Derivation.}
The median of the calm-period energy distribution is the natural temperature because:
\begin{enumerate}
\item At $E = \theta$, $\gamma = 1/2$ — half intervention. During normal operation ($E \approx \theta$ by construction), the system applies moderate damping.
\item The median is robust to outlier spikes in the calibration window. The mean would overstate $\theta$.
\item Setting $\theta = E_\text{median}$ means the system is at its maximum-entropy operating point on average during the calm regime — it neither over-intervenes nor under-intervenes.
\end{enumerate}

\textbf{Alternative derivations of $\theta$:}

Route 1 (Maximum entropy): set $\theta = e^*$ (critical manifold energy) so $\gamma = 1/2$ at stability boundary.

Route 2 (Admissibility): choose $\theta > M_\mathrm{max}\sqrt{M_G}/(\sigma\mu_G)$ to guarantee $\Psi^* > M_\mathrm{max}$.

Route 3 (Chi-squared calibration): $e^* = \chi^2_{Nd,0.99}$, set $\theta = e^*/2$ so $\gamma = 1/3$ at the $99\%$ energy quantile.

%=============================================================
\section{The Canonical ODE}
%=============================================================

\subsection{Locked form}

\begin{lockedbox}
\[
\dot X = \Fbase(X) - \gX(X) - \gamma(X)\,\frac{\inner{\Fbase(X)-\gX(X)}{\gX(X)}}{\norm{\gX(X)}^2}\,\gX(X)
\]
\end{lockedbox}

\subsection{Decomposition}

Define the modified force:
\[
F(X) = \Fbase(X) - \gX(X).
\]
The projection coefficient:
\[
\alpha(X) = \frac{\inner{F(X)}{\gX(X)}}{\norm{\gX(X)}^2}.
\]
The ODE becomes:
\[
\dot X = F(X) - \gamma(X)\alpha(X)\gX(X).
\]

\textbf{Three-term decomposition:}
\begin{enumerate}
\item $\Fbase(X)$: domain-specific nominal field
\item $-\gX(X)$: gradient correction forcing $E$ to decrease
\item $-\gamma\alpha\gX$: adaptive projection removing the component of $F$ aligned with $\gX$
\end{enumerate}

\subsection{Force decomposition}

\[
\dot X = \underbrace{F - \langle F,\hat u\rangle\hat u}_{F_\perp \text{ (always free)}}
+ \underbrace{(1-\gamma)\langle F,\hat u\rangle\hat u}_{(1-\gamma)F_\parallel \text{ (damped)}}
\]
where $\hat u = \gX/\norm{\gX}$.

At $E \to 0$: $\gamma \to 0$, $\dot X \to \Fbase$ (free dynamics).
At $E \to \infty$: $\gamma \to 1$, $\dot X \to F_\perp$ (only safe motion).

%=============================================================
\section{Lyapunov Energy Derivative $\dot E$}
%=============================================================

\subsection{Master identity}

\begin{theorem}[Energy derivative, Lemma 7.7 of v4]
Along any solution of the canonical ODE:
\[
\dot E = (1-\gamma)\left[\inner{\gX}{\Fbase} - \norm{\gX}^2\right]
\]
\end{theorem}

\begin{proof}
\begin{align*}
\dot E &= \langle\gX, \dot X\rangle \\
&= \langle\gX, \Fbase - \gX - \gamma\alpha\gX\rangle \\
&= \langle\gX,\Fbase\rangle - \norm{\gX}^2 - \gamma\alpha\norm{\gX}^2 \\
&= \langle\gX,\Fbase\rangle - \norm{\gX}^2 - \gamma\langle F,\gX\rangle \\
&= \langle\gX,\Fbase\rangle - \norm{\gX}^2 - \gamma\langle\Fbase-\gX,\gX\rangle \\
&= \langle\gX,\Fbase\rangle - \norm{\gX}^2 - \gamma\langle\Fbase,\gX\rangle + \gamma\norm{\gX}^2 \\
&= (1-\gamma)\langle\gX,\Fbase\rangle - (1-\gamma)\norm{\gX}^2 \\
&= (1-\gamma)\left[\langle\gX,\Fbase\rangle - \norm{\gX}^2\right]. \quad\square
\end{align*}
\end{proof}

\subsection{Special cases}

\textbf{Pure geometric flow} ($\Fbase = 0$):
\[
\dot E = -(1-\gamma)\norm{\gX}^2 = -(1-\gamma)\cdot\mfls^2 \leq 0.
\]
Strictly negative whenever $\gX \neq 0$ and $\gamma < 1$.

\textbf{Descent condition} (general $\Fbase$):
\[
\dot E \leq 0 \iff \langle\gX,\Fbase\rangle \leq \norm{\gX}^2.
\]

\textbf{At the maximum-entropy point} ($E = \theta$, $\gamma = 1/2$):
\[
\dot E = \half\left[\langle\gX,\Fbase\rangle - \norm{\gX}^2\right].
\]

%=============================================================
\section{Effective Decay Rate $\rho_\mathrm{eff}$}
%=============================================================

\subsection{Definition}

\begin{lockedbox}
\[
\rho_\mathrm{eff}(t) = -\frac{\dot E(t)}{E(t)}
\]
\end{lockedbox}

\subsection{Derivation in canonical terms}

From the energy derivative:
\[
\rho_\mathrm{eff} = \frac{(1-\gamma)\left[\norm{\gX}^2 - \langle\gX,\Fbase\rangle\right]}{E}.
\]

Under $\Fbase = 0$:
\[
\rho_\mathrm{eff} = \frac{(1-\gamma)\norm{\gX}^2}{E} = \frac{(1-\gamma)\cdot\mfls^2}{E}.
\]

\subsection{Connection to $\sigma$}

From the coercivity bound $\norm{\gX}^2 \geq C_1^2 E$:
\[
\rho_\mathrm{eff} \geq (1-\gamma)C_1^2 = (1-\gamma)\frac{4\sigma^2\mu_G^2}{M_G} = \rho_\text{theoretical}.
\]

The theoretical rate $\rho$ is the worst-case lower bound.
The effective rate $\rho_\mathrm{eff}$ is the instantaneous observed rate.
The ratio $\rho_\mathrm{eff}/\rho \geq 1$ measures how much better the system is
performing relative to the theoretical guarantee.

\subsection{Half-life prediction}

The local exponential half-life is:
\[
\tau_{1/2}(t) = \frac{\ln 2}{\rho_\mathrm{eff}(t)}.
\]
This is the time at which $E$ will reach $E(t)/2$ at the current decay rate.

%=============================================================
\section{Effective Singular Value $\sigma_\mathrm{eff}$}
%=============================================================

\subsection{Abstract definition}

$\sigma = \sigma_{\min}(J(X))$ from Theorem 9.3 of v4 is the abstract worst-case bound.

\subsection{Derivation in canonical terms}

\begin{lockedbox}
\[
\sigma_\mathrm{eff}(t) = \frac{\norm{\gX(t)}\sqrt{M_G}}{2\mu_G\sqrt{E(t)}} = \frac{\mfls_t\sqrt{M_G}}{4\mu_G\sqrt{E_t}}
\]
\end{lockedbox}

\begin{proof}
From the coercivity bound:
\[
\norm{\gX}^2 \geq \frac{4\sigma^2\mu_G^2}{M_G}E.
\]
Treat this as an equality to define the effective $\sigma$ that is tight at the current state:
\[
\norm{\gX}^2 = \frac{4\sigma_\mathrm{eff}^2\mu_G^2}{M_G}E
\implies
\sigma_\mathrm{eff}^2 = \frac{\norm{\gX}^2 M_G}{4\mu_G^2 E}
\implies
\sigma_\mathrm{eff} = \frac{\norm{\gX}\sqrt{M_G}}{2\mu_G\sqrt{E}}.
\]
Since $\mfls = \norm{\gX}$ (Section~4):
\[
\sigma_\mathrm{eff} = \frac{\mfls\sqrt{M_G}}{2\mu_G\sqrt{E}}.
\]
\end{proof}

\subsection{Properties}

$\sigma_\mathrm{eff}(t) \geq \sigma$ always (effective is at least as good as worst case).

At the coercivity lower bound: $\sigma_\mathrm{eff} = \sigma$ (bound is tight when $J$ is at its worst orientation).

$\sigma_\mathrm{eff}$ is a time-varying signal measuring the current Jacobian rank quality — computable at each step.

%=============================================================
\section{Theoretical Rate $\rho$}
%=============================================================

\subsection{Definition}

\begin{lockedbox}
\[
\rho = \frac{4\sigma^2\mu_G^2}{M_G}(1-\gamma_\mathrm{max}), \qquad
\gamma_\mathrm{max} = \frac{E(X_0)}{E(X_0)+\theta}
\]
\end{lockedbox}

\subsection{Derivation}

From Theorem 9.3 of v4 with $\Fbase = 0$:
\[
\dot E = -(1-\gamma)\norm{\gX}^2.
\]
Apply the coercivity lower bound:
\[
\norm{\gX}^2 \geq \frac{4\sigma^2\mu_G^2}{M_G}E.
\]
Since $\gamma \leq \gamma_\mathrm{max} = E_0/(E_0+\theta) < 1$ (monotone in $E$, decreasing, so maximum at $t=0$):
\[
\dot E \leq -(1-\gamma_\mathrm{max})\frac{4\sigma^2\mu_G^2}{M_G}E = -\rho E.
\]
By Grönwall: $E(t) \leq E_0 e^{-\rho t}$.

\subsection{Expression via $\sigma_\mathrm{eff}$}

Substituting $\sigma = \sigma_\mathrm{eff}$:
\[
\rho_\mathrm{eff} = \frac{4\sigma_\mathrm{eff}^2\mu_G^2}{M_G}(1-\gamma)
= \frac{4\cdot\frac{\norm{\gX}^2 M_G}{4\mu_G^2 E}\cdot\mu_G^2}{M_G}(1-\gamma)
= \frac{\norm{\gX}^2}{E}(1-\gamma). \checkmark
\]
Consistent with the direct definition of $\rho_\mathrm{eff}$.

%=============================================================
\section{Admissibility Threshold $\Psi^*$}
%=============================================================

\subsection{Abstract form}

From Theorem 9.4 of v4:
\[
\Psi^* = \frac{\sigma\theta\mu_G}{\sqrt{M_G}}.
\]
This is unimodal with maximum at $R^* = \theta$.

\subsection{Derivation in canonical terms}

\begin{lockedbox}
\[
\Psi^* = \frac{\theta\norm{\gX}}{4\sqrt{E}} = \frac{\theta\cdot\mfls}{4\sqrt{E}}
\]
\end{lockedbox}

\begin{proof}
Substitute $\sigma = \sigma_\mathrm{eff} = \norm{\gX}\sqrt{M_G}/(2\mu_G\sqrt{E})$:
\[
\Psi^* = \frac{\sigma_\mathrm{eff}\theta\mu_G}{\sqrt{M_G}}
= \frac{\norm{\gX}\sqrt{M_G}}{2\mu_G\sqrt{E}}\cdot\frac{\theta\mu_G}{\sqrt{M_G}}
= \frac{\theta\norm{\gX}}{2\sqrt{E}}.
\]
Since $\mfls = \norm{\gX}$:
\[
\Psi^* = \frac{\theta\cdot\mfls}{2\sqrt{E}}.
\]
\end{proof}

\textbf{Note on the factor.} The formula above uses $\sigma_\mathrm{eff}$ (instantaneous effective singular value). Using the abstract worst-case $\sigma$ gives $\Psi^*_\text{abstract} = \sigma\theta\mu_G/\sqrt{M_G}$, while the canonical instantaneous formula gives $\Psi^*(t) = \theta\mfls_t/(2\sqrt{E_t})$. Both are valid; the canonical form is computable at each step.

\subsection{BSDT instantiation}

With $G = \Sigma_0^{-1}$, $\mu_G = 1/\lambda_\mathrm{max}(\Sigma_0)$, $M_G = 1/\lambda_\mathrm{min}(\Sigma_0)$:
\[
\Psi^*_\bsdt = \frac{\sigma\theta/\lambda_\mathrm{max}(\Sigma_0)}{\sqrt{1/\lambda_\mathrm{min}(\Sigma_0)}}
= \frac{\sigma\theta\sqrt{\lambda_\mathrm{min}(\Sigma_0)}}{\lambda_\mathrm{max}(\Sigma_0)}
= \frac{\sigma\theta}{\sqrt{\kappa(\Sigma_0)}\cdot\lambda_\mathrm{max}^{1/2}(\Sigma_0)}.
\]
The canonical form $\theta\mfls/(2\sqrt{E})$ is metric-free — the $\Sigma_0$ terms cancel. $\square$

\subsection{Sensitivity to calibration conditioning}

From the BSDT formula:
\[
\Psi^*_\bsdt \propto \frac{1}{\sqrt{\kappa(\Sigma_0)}}.
\]
A $10\times$ increase in calibration conditioning shrinks $\Psi^*$ by $\sqrt{10} \approx 3.16\times$.
This is the gap closure result G7.2.

%=============================================================
\section{Admissibility Condition}
%=============================================================

\subsection{Statement}

\begin{lockedbox}
\[
\Psi^* > M_\mathrm{max} \iff \frac{\theta\cdot\mfls_t}{2\sqrt{E_t}} > M_\mathrm{max}
\iff \frac{\mfls_t}{\sqrt{E_t}} > \frac{2M_\mathrm{max}}{\theta}
\]
\end{lockedbox}

\subsection{Derivation}

Starting from the abstract condition:
\[
\Psi^* > M_\mathrm{max}
\]
Substituting the canonical form:
\[
\frac{\theta\cdot\mfls}{2\sqrt{E}} > M_\mathrm{max}
\]
Rearranging:
\[
\mfls > \frac{2M_\mathrm{max}\sqrt{E}}{\theta}
\]
Dividing both sides by $\sqrt{E}$:
\[
\frac{\mfls}{\sqrt{E}} > \frac{2M_\mathrm{max}}{\theta}.
\]

\subsection{The canonical admissibility ratio}

\begin{lockedbox}
\[
\mathcal{A}(t) = \frac{\mfls_t}{\sqrt{E_t}} = \frac{\norm{\gX(t)}}{\sqrt{E(t)}}
\]
\end{lockedbox}

This ratio is the canonical admissibility signal. It measures gradient strength relative to energy level.

The admissibility condition is then simply:
\[
\mathcal{A}(t) > \frac{2M_\mathrm{max}}{\theta}.
\]

Connection to $\rho_\mathrm{eff}$:
\[
\rho_\mathrm{eff}(t) = \mathcal{A}(t)^2(1-\gamma(t)).
\]
A system decaying fast is a system that is safe. Fast decay implies admissibility.

%=============================================================
\section{Safety Ratio}
%=============================================================

\subsection{Definition}

\begin{lockedbox}
\[
\mathrm{SafetyRatio}(t) = \frac{\Psi^*(t)}{M_\mathrm{max}} = \frac{\theta\cdot\mfls_t}{2M_\mathrm{max}\sqrt{E_t}}
\]
\end{lockedbox}

\subsection{Derivation}

Divide the admissibility condition by $M_\mathrm{max}$:
\[
\mathrm{SafetyRatio} = \frac{\Psi^*}{M_\mathrm{max}} = \frac{\theta\mfls}{2M_\mathrm{max}\sqrt{E}}.
\]

Substituting $\mfls = \norm{\gX}$, $E$, $\theta$ — all already computed in the canonical loop:

\textbf{Zero additional computation.}

\subsection{Regimes}

\begin{center}
\begin{tabular}{lll}
\toprule
SafetyRatio & Regime & Interpretation \\
\midrule
$> 1$ & Admissible & $\Psi^* > M_\mathrm{max}$; ultimate boundedness guaranteed \\
$= 1$ & Boundary & Exactly at the admissibility threshold \\
$< 1$ & Inadmissible & Disturbance can overcome canonical brake \\
$\to 0$ & Critical & $\mfls \to 0$ or $E \to \infty$; system losing gradient signal \\
\bottomrule
\end{tabular}
\end{center}

\subsection{Connection to effective decay rate}

\[
\mathrm{SafetyRatio}(t) = \frac{\theta}{2M_\mathrm{max}}\mathcal{A}(t) = \frac{\theta}{2M_\mathrm{max}}\sqrt{\frac{\rho_\mathrm{eff}(t)}{1-\gamma(t)}}.
\]

This connects the safety guarantee directly to the observable performance metric.

%=============================================================
\section{Alignment Angle $\cos\theta_t$}
%=============================================================

\subsection{Definition}

\begin{lockedbox}
\[
\cos\theta_t = \frac{\inner{\Fbase(X_t)}{\gX(X_t)}}{\norm{\Fbase(X_t)}\cdot\norm{\gX(X_t)}} \in [-1,+1]
\]
\end{lockedbox}

\subsection{Derivation from energy descent}

The descent condition $\dot E \leq 0$ requires:
\[
\inner{\gX}{\Fbase} \leq \norm{\gX}^2.
\]
Writing $\inner{\gX}{\Fbase} = \norm{\gX}\norm{\Fbase}\cos\theta_t$:
\[
\norm{\Fbase}\cos\theta_t \leq \norm{\gX}.
\]
When $\cos\theta_t \leq 0$: descent is automatic regardless of $\norm{\Fbase}$.
When $\cos\theta_t > 0$: descent requires $\norm{\Fbase}\cos\theta_t < \norm{\gX}$,
i.e.\ $\Fbase$ is not pushing into the energy gradient too strongly.

\subsection{Energy derivative in terms of $\cos\theta_t$}

\[
\dot E = (1-\gamma)\norm{\gX}\left[\norm{\Fbase}\cos\theta_t - \norm{\gX}\right].
\]

\subsection{Geometric interpretation}

\begin{center}
\begin{tabular}{ll}
\toprule
Value & Meaning \\
\midrule
$\cos\theta_t \to +1$ & $\Fbase$ co-aligned with $\gX$ — coherent collapse pressure \\
$\cos\theta_t \approx 0$ & $\Fbase$ tangent to energy level set — no radial progress \\
$\cos\theta_t \to -1$ & Anti-aligned — restoring competition, $E$ falls without $\Fbase$ help \\
\bottomrule
\end{tabular}
\end{center}

\subsection{Collapse angle $\tan\theta_t$}

From §XIV of the anatomy document:
\[
\tan\theta_t = \frac{\norm{F_\perp}}{\norm{F_\parallel}} = \sqrt{\frac{1}{\cos^2\theta_t} - 1} = \frac{\sqrt{1-\cos^2\theta_t}}{\cos\theta_t}.
\]

$\tan\theta_t \to 0$: trajectory fully collapse-directed.
$\tan\theta_t \to \infty$: motion fully perpendicular, locally safe.
$\tan\theta_t = 1$ ($\theta_t = 45°$): critical balance.

%=============================================================
\section{Misalignment Angle $\psi_t$ (Sixth Signal)}
%=============================================================

\subsection{Definition}

\begin{lockedbox}
\[
\cos\psi_t = \frac{\inner{\tilde G_t}{g_t}_F}{\norm{\tilde G_t}_F\cdot\norm{g_t}}
\]
\end{lockedbox}

where $g_t = \nabla_S E_\mathrm{BS}(S_t) \in \R^4$ is the channel-space gradient
and $\tilde G_t = J_t\T g_t \in \R^{Nd}$ is its state-space pullback.

\subsection{Derivation}

The BSDT master operator uses $g_t$ (channel space) in the projection.
The canonical ODE uses $\tilde G_t$ (state-space pullback).

These differ by the Jacobian map: $\tilde G_t = J_t\T g_t$.

The angle between them:
\[
\cos\psi_t = \frac{\inner{\tilde G_t, g_t}_F}{\norm{\tilde G_t}_F\norm{g_t}}
= \frac{R_t}{\norm{\tilde G_t}_F\norm{g_t}}
\]
where $R_t = \sum_k [g_t]_k\langle\partial_X\delta_k, g_t\rangle_F$ is the control coupling sum.

\subsection{What $\psi_t$ measures}

$\psi_t \approx 0$: channel-space risk and physical dynamics pointing in the same direction.
Collapse is real. All five prior signals are reliable.

$\psi_t \approx \pi/2$: complete misalignment.
Abstract risk elevated but physical dynamics not aligned with collapse.
This is a false alarm.

\subsection{False alarm resolution}

Previously: high MFLS but no actual collapse. Unexplained by any single-space signal.

Now: high channel MFLS ($\norm{g_t}$ large) with $\psi_t$ near $\pi/2$
means $\tilde G_t \perp g_t$ — physical dynamics are tangential to the collapse direction.
The alarm is real in channel space but not in state space.
$\psi_t$ discriminates between the two.

\subsection{Performance gap formula}

The performance difference between BSDT and CEK-v4:
\[
\Delta\text{performance}(t) \approx f(\psi_t) + \kappa(\mathcal{R}_t)^{-1}
\]
where:
\begin{itemize}
\item $\psi_t$ large: BSDT projects using wrong direction, canonical corrects
\item $\kappa(\mathcal{R}_t)$ large: channels nearly collinear, canonical compression is more robust
\end{itemize}

%=============================================================
\section{Channel Recovery Operator $\mathcal{R}$}
%=============================================================

\subsection{The recovery formula}

\begin{lockedbox}
\[
S_t = \mathcal{R}(X_t)\,\gX(t), \qquad \mathcal{R}(X) = \frac{1}{2}A^{-1}(JJ\T)^{-1}J \in \R^{k\times n}
\]
\end{lockedbox}

\subsection{Derivation}

From $\gX = 2J\T GS$, left-multiply by $J$:
\[
J\gX = 2JJ\T GS.
\]
Since $JJ\T \in \R^{k\times k}$ is invertible under the rank condition:
\[
GS = \frac{1}{2}(JJ\T)^{-1}J\gX.
\]
Since $G = A \succ 0$ is invertible:
\[
S = \frac{1}{2}G^{-1}(JJ\T)^{-1}J\gX = \frac{1}{2}A^{-1}(JJ\T)^{-1}J\gX = \mathcal{R}\gX.
\]

\subsection{Per-channel recovery}

\[
[\delta_k] = [S]_k = [\mathcal{R}\gX]_k, \quad k \in \{C,G,A,T\}.
\]

\subsection{Attribution}

\[
a_k(t) = \frac{\tilde S_k(t)^2}{\norm{\tilde S_t}^2}, \quad \tilde S_k = \frac{S_k}{\mu_k^\text{normal}}, \quad \sum_k a_k = 1.
\]

\subsection{Stability of recovery}

\begin{proposition}[Lipschitz stability]
For perturbed gradient $\tilde\gX = \gX + \varepsilon$ with $\norm{\varepsilon} \leq \delta$:
\[
\norm{\tilde S - S} \leq \norm{\mathcal{R}}_\mathrm{op}\cdot\delta = \frac{\delta}{2\sigma_{\min}(A)\cdot\sigma_{\min}(JJ\T)^{1/2}\cdot\sigma_{\min}(J)}.
\]
\end{proposition}

\begin{proof}
$\tilde S - S = \mathcal{R}\varepsilon$, so $\norm{\tilde S - S} \leq \norm{\mathcal{R}}_\mathrm{op}\norm{\varepsilon}$.
Expanding $\norm{\mathcal{R}}_\mathrm{op} \leq \frac{1}{2}\norm{A^{-1}}_\mathrm{op}\norm{(JJ\T)^{-1}}_\mathrm{op}\norm{J}_\mathrm{op}$
and substituting singular value bounds gives the result.
\end{proof}

\subsection{Condition number of recovery}

\[
\kappa(\mathcal{R}) = \frac{\sigma_{\max}(\mathcal{R})}{\sigma_{\min}(\mathcal{R})}.
\]

Recovery is well-conditioned when $\kappa(\mathcal{R}) \ll 1$.

Recovery fails (ill-conditioned) when:
\begin{itemize}
\item Channels nearly collinear: $\sigma_{\min}(JJ\T) \approx 0$
\item Coupling matrix ill-conditioned: $\kappa(A) \gg 1$
\item Jacobian rank-deficient: $\sigma_{\min}(J) \approx 0$ (channel has zero sensitivity)
\end{itemize}

%=============================================================
\section{Ultimate Boundedness Threshold $\Psi^*_\mathrm{rob}$}
%=============================================================

\subsection{Standard case}

From Theorem 9.4 of v4:
\[
\Psi_\mathrm{std}(R) = (1-\gamma(R))\cdot 2\sigma\sqrt{\frac{\mu_G^2}{M_G}R}
= \frac{\theta}{R+\theta}\cdot C_1\sqrt{R}.
\]

Setting $u = \sqrt{R}$:
\[
\Psi_\mathrm{std}(u) = \frac{\theta C_1 u}{u^2 + \theta}.
\]

Maximum at $u = \sqrt{\theta}$ (i.e.\ $R = \theta$):
\[
\Psi^*_\mathrm{std} = \frac{C_1\sqrt{\theta}}{2}.
\]

\subsection{Robust case with adversarial disturbance}

With two-sided coercivity $C_1\sqrt{E} \leq \norm{\gX} \leq C_2\sqrt{E}$:
\[
\Psi_\mathrm{rob}(R) = \frac{\theta C_1\sqrt{R}}{R + M_\adv C_2\sqrt{R} + \theta} - M_\adv.
\]

Setting $u = \sqrt{R}$:
\[
\frac{d\Psi_\mathrm{rob}}{du} = \frac{\theta C_1(\theta - u^2)}{(u^2 + M_\adv C_2 u + \theta)^2}.
\]

Positive for $u < \sqrt{\theta}$, negative for $u > \sqrt{\theta}$: unique maximum at $u = \sqrt{\theta}$, i.e.\ $R = \theta$.

\textbf{Unimodality proved.}

\subsection{Admissibility condition (robust)}

\[
M_\nom < \Psi^*_\mathrm{rob} = \Psi_\mathrm{rob}(\theta) = \frac{\theta C_1\sqrt{\theta}}{\theta + M_\adv C_2\sqrt{\theta} + \theta} - M_\adv.
\]

%=============================================================
\section{The Complete Dashboard}
%=============================================================

\subsection{All signals from two objects}

Every signal in the canonical dashboard is a function of $\gX$ and $E$:

\begin{center}
\renewcommand{\arraystretch}{1.4}
\begin{tabular}{llll}
\toprule
Signal & Formula & Objects & Section \\
\midrule
$\mfls_t$ & $\norm{\gX}$ & $\gX$ & \S4 \\
$E_t$ & $S\T GS$ & $S, G$ & \S2 \\
$\gamma_t$ & $E/(E+\theta)$ & $E,\theta$ & \S5 \\
$\dot E_t$ & $(1-\gamma)(\inner{\gX}{\Fbase}-\norm{\gX}^2)$ & $\gX,E,\theta,\Fbase$ & \S7 \\
$\rho_\mathrm{eff}(t)$ & $\norm{\gX}^2(1-\gamma)/E$ & $\gX,E,\theta$ & \S8 \\
$\cos\theta_t$ & $\inner{\Fbase,\gX}/(\norm{\Fbase}\norm{\gX})$ & $\gX,\Fbase$ & \S11 \\
$\tan\theta_t$ & $\sqrt{1/\cos^2\theta_t - 1}$ & $\cos\theta_t$ & \S11 \\
$\sigma_\mathrm{eff}(t)$ & $\norm{\gX}\sqrt{M_G}/(2\mu_G\sqrt{E})$ & $\gX,E,\mu_G,M_G$ & \S9 \\
$\Psi^*(t)$ & $\theta\norm{\gX}/(2\sqrt{E})$ & $\gX,E,\theta$ & \S10 \\
$\mathcal{A}(t)$ & $\norm{\gX}/\sqrt{E}$ & $\gX,E$ & \S10 \\
$\mathrm{SafetyRatio}(t)$ & $\theta\norm{\gX}/(2M_\mathrm{max}\sqrt{E})$ & $\gX,E,\theta$ & \S10 \\
$\tau_{1/2}(t)$ & $\ln 2/\rho_\mathrm{eff}(t)$ & $\rho_\mathrm{eff}$ & \S8 \\
$\psi_t$ & $\arccos(R_t/(\norm{\tilde G_t}\norm{g_t}))$ & $J,g_t,\gX$ & \S12 \\
$S_t$ (channels) & $\mathcal{R}(X_t)\gX(t)$ & $J,A,\gX$ & \S13 \\
$a_k(t)$ & $\tilde S_k^2/\norm{\tilde S}^2$ & $S_t$ & \S13 \\
\bottomrule
\end{tabular}
\end{center}

\subsection{Dependency graph}

\begin{verbatim}
X
├─ S = B(X)         ──► E = S^T G S
│                       │
│                       ├─ γ = E/(E+θ)
│                       ├─ γ_max = E_0/(E_0+θ)
│                       └─ ρ = 4σ²μ²_G/M_G · (1-γ_max)
│
└─ J = ∂S/∂X        ──► g_X = 2J^T G S = MFLS
                         │
                         ├─ MFLS = ||g_X||
                         ├─ σ_eff = MFLS·√M_G / (2μ_G·√E)
                         ├─ Ψ* = θ·MFLS / (2√E)
                         ├─ SafetyRatio = Ψ*/M_max
                         ├─ ρ_eff = MFLS²·(1-γ)/E
                         ├─ cos θ_t = <F_base, g_X> / (||F_base||·MFLS)
                         └─ R = R(X)·g_X → channels S_t → a_k(t)
\end{verbatim}

\subsection{Computational algorithm (per step $t$)}

\begin{enumerate}
\item Compute $\tilde X_t = X_t - \mathbf{1}\mu_0\T$
\item Compute $S_t = B(X_t) = (\delta_C,\delta_G,\delta_A,\delta_T)\T$
\item Compute $J_t = \partial B/\partial X|_{X_t}$
\item Compute $\gX(t) = 2J_t\T G S_t$ \hfill \textit{(canonical gradient)}
\item Compute $E_t = S_t\T G S_t$ \hfill \textit{(energy)}
\item Compute $\gamma_t = E_t/(E_t+\theta)$ \hfill \textit{(gain)}
\item Compute $\mfls_t = \norm{\gX(t)}$ \hfill \textit{(free from step 4)}
\item Compute $\rho_\mathrm{eff}(t) = \mfls_t^2(1-\gamma_t)/E_t$
\item Compute $\Psi^*(t) = \theta\mfls_t/(2\sqrt{E_t})$
\item Compute $\mathrm{SafetyRatio}(t) = \Psi^*(t)/M_\mathrm{max}$
\item Compute $\cos\theta_t = \inner{F_\mathrm{base}(X_t)}{\gX(t)}/(\norm{F_\mathrm{base}}\mfls_t)$
\item Compute $\mathcal{R}_t = \frac{1}{2}A^{-1}(J_tJ_t\T)^{-1}J_t$ \hfill \textit{($4\times4$ inversion)}
\item Compute $S_t = \mathcal{R}_t\gX(t)$ \hfill \textit{(channel recovery)}
\item Compute $a_k(t) = \tilde S_k(t)^2/\norm{\tilde S_t}^2$ \hfill \textit{(attribution)}
\item Compute $\psi_t = \arccos(R_t/(\norm{\tilde G_t}_F\norm{g_t}))$
\item Update ODE: $\dot X = \Fbase - \gX - \gamma_t\alpha_t\gX$
\end{enumerate}

%=============================================================
\section{BSDT--Canonical Juxtaposition}
%=============================================================

\subsection{Identity table}

\begin{center}
\renewcommand{\arraystretch}{1.3}
\begin{tabularx}{\textwidth}{l X X l}
\toprule
Object & BSDT form & Canonical form & Equal? \\
\midrule
State & $X \in \R^{N\times d}$ & $X \in \R^{N\times d}$ & Identical \\
Energy & $e_t = \norm{\tilde X\Sigma_0^{-1/2}}_F^2$ & $E = S\T GS$ & Equal under $G=A$ \\
Gradient & $\tilde G_t = \sum_k[AS]_k\partial_X\delta_k$ & $\gX = 2J\T GS$ & Identical \\
Gain & $\gamma^* = e_t/(e_t+\theta)$ & $\gamma = E/(E+\theta)$ & Identical \\
ODE & $\dot X = F_t - \gamma^*\frac{\inner{F_t,g_t}}{\norm{g_t}^2}g_t$ & locked canonical ODE & Identical \\
MFLS & $2\norm{\tilde X\Sigma_0^{-1}}_F$ & $\norm{\gX}$ & Identical \\
Rate $\rho$ & not proved & $4\sigma^2\mu_G^2/M_G(1-\gamma_\mathrm{max})$ & canonical proves BSDT \\
$\Psi^*$ & not proved & $\theta\mfls/(2\sqrt{E})$ & canonical proves BSDT \\
Channels & $(\delta_C,\delta_G,\delta_A,\delta_T)$ explicit & recovered via $\mathcal{R}\gX$ & $4\times4$ inversion \\
$\psi_t$ & not in BSDT & $\arccos(R_t/\norm{\tilde G_t}\norm{g_t})$ & new in canonical \\
Extensions & none & 7 extensions & canonical adds \\
\bottomrule
\end{tabularx}
\end{center}

\subsection{The master identity}

\[
\underbrace{S_t}_{\text{BSDT channels}}
=
\underbrace{\frac{1}{2}A^{-1}(J_tJ_t\T)^{-1}J_t}_{\text{recovery operator }\mathcal{R}(X_t)}
\underbrace{\gX(t)}_{\text{canonical gradient}}
\]

BSDT is on the left. Canonical is on the right.
The recovery operator in the middle is the bridge.
The equation says they are the same object, read from different directions.

\subsection{Performance gap}

BSDT uses $g_t$ (channel space) in the projection.
Canonical uses $\tilde G_t$ (state-space pullback).
The gap is $\psi_t$:

\[
\Delta\text{performance}(t) \sim \psi_t + \kappa(\mathcal{R}_t)^{-1}.
\]

The v34$\to$v58 empirical progression was systematically reducing $\psi_t$
and $\kappa(\mathcal{R}_t)$ without knowing it:

\begin{center}
\begin{tabular}{lll}
\toprule
Version & Implicit improvement & Canonical equivalent \\
\midrule
v34 & Base & No channel recovery \\
v48 & 4-quadrant modulator & Partial recovery via $(G_\mathrm{mem},T_\mathrm{mem})$ \\
v52 & Clip structure & Excludes high $\kappa(\mathcal{R})$ regimes \\
v55 & State-dependent $\theta_G(t)$ & Ricci-flow metric adaptation \\
v58 & Asymmetric clamp & Ill-conditioned recovery region excluded \\
\bottomrule
\end{tabular}
\end{center}

%=============================================================
\section{Summary Reference}
%=============================================================

\subsection{Every formula}

\begin{resultbox}
\begin{align*}
S &= B(X) = (\delta_C,\delta_G,\delta_A,\delta_T)\T &&\text{feature vector} \\[4pt]
E &= S\T GS &&\text{Lyapunov energy} \\[4pt]
J &= \partial S/\partial X \in \R^{k\times n} &&\text{Jacobian} \\[4pt]
\gX &= 2J\T GS &&\text{canonical gradient} \\[4pt]
\mfls &= \norm{\gX} = 2\norm{J\T GS} &&\text{field line score} \\[4pt]
\gamma &= \frac{E}{E+\theta} &&\text{adaptive gain} \\[4pt]
\theta &= \mathrm{median}_{t\in\text{calm}}\{E(t)\} &&\text{intervention cost} \\[4pt]
\dot E &= (1-\gamma)\left[\inner{\gX}{\Fbase}-\norm{\gX}^2\right] &&\text{energy derivative} \\[4pt]
\rho_\mathrm{eff} &= \frac{\norm{\gX}^2(1-\gamma)}{E} = \frac{\mfls^2(1-\gamma)}{E} &&\text{instantaneous rate} \\[4pt]
\sigma_\mathrm{eff} &= \frac{\mfls\sqrt{M_G}}{2\mu_G\sqrt{E}} &&\text{effective singular value} \\[4pt]
\rho &= \frac{4\sigma^2\mu_G^2}{M_G}(1-\gamma_\mathrm{max}) &&\text{theoretical rate} \\[4pt]
\Psi^* &= \frac{\theta\cdot\mfls}{2\sqrt{E}} &&\text{admissibility threshold} \\[4pt]
\mathcal{A} &= \frac{\mfls}{\sqrt{E}} = \frac{\norm{\gX}}{\sqrt{E}} &&\text{admissibility ratio} \\[4pt]
\text{SafetyRatio} &= \frac{\theta\cdot\mfls}{2M_\mathrm{max}\sqrt{E}} &&\text{real-time safety signal} \\[4pt]
\cos\theta_t &= \frac{\inner{\Fbase}{\gX}}{\norm{\Fbase}\norm{\gX}} &&\text{alignment} \\[4pt]
\cos\psi_t &= \frac{\inner{\tilde G_t}{g_t}_F}{\norm{\tilde G_t}_F\norm{g_t}} &&\text{misalignment (sixth signal)} \\[4pt]
\mathcal{R} &= \frac{1}{2}A^{-1}(JJ\T)^{-1}J &&\text{channel recovery operator} \\[4pt]
S_t &= \mathcal{R}(X_t)\gX(t) &&\text{channel recovery formula} \\[4pt]
a_k &= \tilde S_k^2/\norm{\tilde S}^2 &&\text{channel attribution}
\end{align*}
\end{resultbox}

\subsection{The two-object reduction}

\textbf{Everything derives from $\gX$ and $E$.}

$\gX$ carries: MFLS, $\sigma_\mathrm{eff}$, $\Psi^*$, SafetyRatio, $\cos\theta_t$, channel recovery.

$E$ carries: $\gamma$, $\rho_\mathrm{eff}$, $\tau_{1/2}$, admissibility condition.

The entire framework is a set of scalar and matrix operations on two vectors computed once per step.

\end{document}
